\begin{itemize}
    \item \textbf{Number of subjects to be accrued/evaluated: } 
% e.g. 600 (300 per arm)
% List number of subjects planned and expected to contribute to the primary analysis. 
% If estimation was based on medical/regulatory reasons, 
% state so and delete all bullet points in this section below.

    \item \textbf{Type-I error rate (one or two-sided): } % e.g. one-sided 0.025, two-sided 0.05
    \item \textbf{Power: } % e.g. 80\%, 90\%, power achieved with planned sample size
    \item \textbf{Assumed values in the estimation: } % may comment the following bulelts if not applicable 
    \begin{itemize}
        \item \textbf{Treatment effect size: }
% Describe expected or desired/clinically meaningful treatment effect size, 
% e.g., HR= 0.60 (median xx vs. yy months), 
% or difference of rate of 0.30 (rate of 0.10 vs. 0.40). 
% For non-inferiority or equivalence designs, e.g. treatment group difference=0. 
% Include specific rates for the treatment and control groups if applicable.
        \item \textbf{Standard deviation: } % List value if the outcome is a continuous variable. Comment this otherwise. 
        \item \textbf{Dropout/missing rate: } % Expected missing data rate, e.g. 10\%
        \item \textbf{Others: } % e.g. over-dispersion rate in a negative-binomial model, projected accrual rate. Comment this otherwise.
    \end{itemize}
    \item \textbf{Rationale provided for assumed values: } % Include any rationale for assumptions described above to size the study.
    \item \textbf{Number of events required to stop/complete the trial: } % List the number of events required if endpoint is a time-to-event endpoint
    \item \textbf{Calculation verification: }
% Yes or No
% If no, need a comment to the sponsor; 
% if the difference between the reviewer’s and the sponsor’s is small, may answer ‘yes’. 
% Briefly describe software/procedure used in the verification, 
% e.g., PASS 13 two-sample t-test, simulation in SAS using PROC PHREG. 
% If sample size is driven by safety, provide justification/explanation for relevant parameters considered.
\end{itemize}